\section{Tokenization}

\textbf{Tokenization} là quá trình chuyển đổi văn bản thô thành các đơn vị nhỏ hơn (tokens) mà máy tính có thể xử lý. Quá trình tokenization được thực hiện thông qua 'TfidfVectorizer' với các tham số:

\begin{itemize}
    \item \textbf{Token pattern:} \texttt{r'\textbackslash b\textbackslash w+\textbackslash b'} - chỉ giữ lại các từ hoàn chỉnh
    \item \textbf{N-gram range:} (1, 2) - tạo unigrams và bigrams để capture cả từ đơn và cụm từ
    \item \textbf{Min/Max document frequency:} Loại bỏ từ xuất hiện quá ít (<2 documents) hoặc quá nhiều (>95\% documents)
    \item \textbf{Max features:} 25,000 từ quan trọng nhất được giữ lại
\end{itemize}

\noindent Việc tokenization cần thiết vì: Chuẩn hóa input, chuyển văn bản liên tục thành format structured. Xây dựng từ điển có kiểm soát từ training data. Chỉ giữ lại những từ/cụm từ có giá trị phân biệt

\section{Vectorization}

\textbf{Vectorization} chuyển đổi tokens thành vectors số học mà machine learning models có thể xử lý. Sử dụng \textit{TF-IDF (Term Frequency-Inverse Document Frequency)} với ưu điểm:

\begin{itemize}
    \item \textbf{Trọng số thông minh}: Từ xuất hiện nhiều trong document nhưng ít trong corpus sẽ có trọng số cao
    \item \textbf{Giảm ảnh hưởng stopwords}: Các từ phổ biến được tự động giảm trọng số
    \item \textbf{Normalize length}: Các document khác độ dài được chuẩn hóa về cùng scale
\end{itemize}

\noindent\textbf{Kết hợp features:}
Text features được \textit{concatenate} với 30 numerical features (statistical + subject categories) tạo thành final feature matrix \texttt{(samples, 25,030)}. Việc kết hợp này cho phép model tận dụng cả \textit{semantic information} từ text và \textit{statistical patterns} từ numerical features.

\noindent\textbf{Tầm quan trọng cho ML/DL models:}
\begin{itemize}
    \item \textbf{Machine learning models} chỉ hiểu số, không hiểu text
    \item \textbf{Consistency}: Đảm bảo mọi input có cùng dimension và format
    \item \textbf{Numerical stability}: TF-IDF tạo ra values trong khoảng ổn định cho optimization algorithms
    \item \textbf{Sparse representation}: Tiết kiệm memory với sparse matrices (non-zero values)
    \item \textbf{Scalability}: Cho phép xử lý large vocabulary một cách hiệu quả
\end{itemize}

\noindent Quan trọng nhất, vectorizer được \textbf{fit chỉ trên training data} và \textit{transform} test data, đảm bảo không có information leakage từ test set về training process.